\section{Guarantees and approximation boundary}
\label{sec:theory}

Results below hold for the finite local mixture in Eq.~\eqref{eq:branch};
they do not certify calibration to the true diffusion posterior.

\begin{proposition}[Mixture covariance factor]
Let $\bar m=\sum_kw_km_k$.  Then
\begin{equation}
 \Cov(x)=LL^\top,\quad
 L=\left[\sqrt{w_k}U_k,\sqrt{w_k}(m_k-\bar m)\right]_{k=1}^K.
\end{equation}
Moreover, $\rank(L)\leq\sum_kr_k+K_+-1$, where $K_+$ is the number of
positive-weight branches.
\end{proposition}
\noindent
This is total covariance.  The mean offsets sum to zero, giving $K_+-1$;
pushing $L$ through an action gives $G_a$ and the same bound.

\begin{theorem}[Information sandwich]
Under independent Gaussian action noise, Eq.~\eqref{eq:mi-decomposition} is
exact for the represented latent $C$.  Its information satisfies
Eq.~\eqref{eq:sandwich}.
\end{theorem}
\noindent
The chain rule gives
$\I(Z,\xi_Z;y)=\I(Z;y)+\I(\xi_Z;y\mid Z)$.  The latter is a weighted
linear-Gaussian log determinant.  Nonnegativity and
$\I(Z;y)\leq H(Z)$ give the lower bound and entropy cap.  A Gaussian maximizes
entropy at fixed covariance, giving the moment-Gaussian upper bound.

\begin{corollary}[Safe one-step screening]
For an equal-cost one-step objective, if $U_a<\max_bL_b$, action $a$ is not
optimal under the mixture surrogate.
\end{corollary}
\noindent
This does not apply to an $H=2$ root: future value may reverse the immediate
ordering.  Confirmatory paired $H=1/H=2$ runs score all roots through the same
estimator; screening remains an optional $H=1$ primitive.

\begin{theorem}[Operator sufficiency]
Let $Q_a$ span $\widetilde G_a$ exactly.  The projected observation $s_a$ in
Eq.~\eqref{eq:projected-observation} is sufficient for the component label and
local coefficient.  Hence it preserves mixture likelihood ratios, mutual
information, and the posterior update, with
\begin{equation}
 \dim(s_a)\leq\sum_kr_k+K_+-1.
\end{equation}
\end{theorem}
\noindent
After whitening, a vector orthogonal to $Q_a$ removes all factors and mean
offsets, leaving standard Gaussian noise independent of $(Z,\xi_Z)$.
Factorization gives sufficiency.  The numerical SVD implementation is
approximate if it drops a nonzero singular value.

\begin{proposition}[No outcome adaptivity for one Gaussian]
For one frozen linear-Gaussian component, independent Gaussian measurements,
and outcome-independent feasibility/costs, an optimal finite-horizon
expected-information policy can be chosen open loop.
\end{proposition}
\noindent
Posterior covariance depends on selected actions, not their values, so
backward induction yields an action-history-only continuation.  In \method{},
adaptive value can arise from updated mode weights/means, relinearization,
operator uncertainty, or outcome-dependent constraints.

\paragraph{Conditional $H=2$ error statement.}
Let $\mathcal B$ be the retained root beam, let
$\Delta_{\mathcal B}=\max_aQ(a)-\max_{a\in\mathcal B}Q(a)$, and suppose every
estimated root, future action, and fantasy expectation is uniformly accurate
within $\varepsilon_0,\varepsilon_1,\varepsilon_F$, respectively.  Then the
selected root $\widehat a$ obeys
\begin{equation}
 \max_aQ(a)-Q(\widehat a)
 \leq\Delta_{\mathcal B}
 +2[\varepsilon_0+\gamma(\varepsilon_1+\varepsilon_F)].
 \label{eq:regret}
\end{equation}
This follows from the 1-Lipschitz maximum and empirical-argmax inequalities.
It is conditional, not a finite-sample certificate: mixture information
density has an unbounded negative tail, nested maxima are biased, and
$\Delta_{\mathcal B}$ is unbounded here.  Monte Carlo errors are descriptive;
cap/beam sensitivity is tested empirically.

\paragraph{Complexity.}
With $T$ reverse steps, $K$ branches, $J$ acquisition rounds, and rank $r$,
central differences use
\begin{equation}
    KT+2KrJ
    \label{eq:nfe}
\end{equation}
image-equivalent denoiser evaluations.  Planning adds zero NFE but nonzero
wall time.  Headline $K=8,r=2,T=32,J=4$ uses 384 NFE, matching a
$K=12,T=32$ point ensemble.  Perturbations and branches are batchable, but we
report calls, launches, memory, and latency.
